\begin{Huge}
    Informatik-FAQ
\end{Huge}

\begin{block}{H\"aufig gestellte Fragen zum Studium}

\question{Lernt man im Studium, wie man programmiert?}{
    Ja, aber auf einer sehr eigenst\"andigen Basis. Man bekommt einen \"Uberblick \"uber die Sprache(n) und \"uber dahinterstehende Programmierkonzepte. Praktische Programmiererfahrung muss man sich aber gr\"o"stenteils selbst erarbeiten.
}

\question{Welche Programmiersprachen macht man da so?}{
    Das ist von Prof zu Prof und von Veranstaltung zu Veranstaltung unterschiedlich. In den ersten beiden Semestern lernt man meistens Racket und Java.
}

\question{Muss man programmieren k\"onnen, um das Studium anzufangen?}{
    Nein. Die Vorlesung beginnt absolut bei 0, um allen den Einstieg zu erm\"oglichen.
}

\question{Muss man gut in Mathe sein?}{
    Man muss kein Mathe-Genie sein, man sollte Mathe aber nicht hassen. Es ist gerade am Anfang viel Mathe.
}

\question{Brauche ich einen eigenen Laptop?}{
    Ist empfehlenswert. Die Anzahl an Rechnern in den Rechnerr\"aumen ist begrenzt und mit dem eigenen Laptop ist man um einiges flexibler. Gerade im ersten Studienjahr braucht man auch nicht unbedingt das leistungsst\"arkste Modell.
}

\question{Wie ist die Frauenquote so?}{
    22,3\%, die Tendenz ist aber leicht steigend. Wir w\"urden uns freuen, wenn die Frauenquote noch weiter steigt.
}

\question{Ich zocke total gerne, hab ich das Zeug, um Informatik zu studieren?}{
    Informatik und Zocken sind erstmal unterschiedliche Dinge. Im Studium musst du analytisches Denken entwickeln, Spa"s am Ausprobieren haben und nicht abgeneigt sein, dich mit Mathe zu besch\"aftigen. Zocken kann einige dieser Eigenschaften f\"ordern, ist aber auch nicht der einzige Weg, diese zu erlangen, und deshalb keine Voraussetzung.
}

\question{Was arbeitet man danach so?}{
    Alle Bereiche der IT-Branche, z. B. Softwareentwicklung und -Beratung, Hardware-Entwicklung, Automatisierung, Automobilindustrie, Unternehmensberatung, Handel, Banken, Versicherungen...
}

\question{Gibt es Praktika?}{
    Im normalen Studienverlauf ist kein berufsorientiertes Praktikum vorgesehen, viele arbeiten aber parallel als Werksstudent*in oder man macht ein Kurzpraktikum in den Semesterferien.
}

\question{Kann man ein Auslandssemster machen?}{
    Klar, geht immer. T\"ubingen nimmt am ERASMUS-Programm teil, die Organisation ist aber langwierig und man sollte sich fr\"uh (ein Jahr vorher) drum k\"ummern.
}

\question{Wie ist da so der NC?}{
    Gibt es keinen, der Studiengang ist zulassungsfrei.
}
\end{block}
